\section{Berechnung der absoluten Feuchte in \texorpdfstring{$\mathrm{g} / \mathrm{kg}$}{g / kg}} 

%\section{Berechnung der absoluten Feuchte in \textorpdfstring{$\text{g} / \text{kg}$}{g / kg}} 


\subsection{Messwerte in der Dämmschichtrocknung}
Berechnung der absoluten Feuchte aus gemessenen Werten:

\begin{itemize}
\item \textbf{relative Feuchtigkeit:} rF in Prozent im Wertebereich $[0 .. 100]$
\item \textbf{Temperatur:} T in °C: 
\item calcValue: absolute Feuchte aF in $\frac{\mathrm{g}}{\mathrm{kg}}$
\end{itemize}

In der POLYGON App und in der Apama-Regel in der CoT wird die Berechnug der absoluten Feuchte in $\frac{\mathrm{g}}{\mathrm{kg}}$ wie folgt durchgeführt:

\begin{minted}[
    bgcolor=gray!10,   % Grauer Hintergrundkasten
    fontsize=\small,   % Angenehme Schriftgröße
    breaklines         % Automatischer Zeilenumbruch
]{java}  
    const calcValue = ((((rF / 100) * (6.112 * Math.pow(2.7182881, (17.62 * T) / (243.12 + T)) * 100)) / 461.5 / (T + 273)) * 1000) / 1.2
\end{minted}

\subsection{Herleitung und Analyse der Formel}
Implementierte Formel:
\begin{equation}
    \text{aF}  
    = \frac{\frac{\frac{\text{rF}}{100} \cdot 6{,}112 \cdot 2{,}7182881^{\frac{17{,}62 \cdot \vartheta}{243{,}12 + \vartheta}} \cdot 100}
                  {461{,}5 \cdot (\vartheta + 273)} \cdot 1000}
            {1{,}2} 
\end{equation}

Schrittweise Vereinfachung:
\begin{equation}
    \text{aF}  
    = \frac{\text{rF} \cdot 6{,}112 \cdot e^{\frac{17{,}62 \cdot \vartheta}{243{,}12 + \vartheta}} }
            {461{,}5 \cdot (\vartheta + 273)}
        \cdot \frac{1000}{1,2}
\end{equation}

Abbilden auf die  \textbf{Tetens-Gleichung} bzw. \textbf{Magnus Gleichung} und anpassen auf die Einheit $\frac{\text{g}}{\text{kg}}$:
\begin{equation}
    \text{aF}  
    = \frac{\text{rF} \cdot 611{,}2\, \frac{\text{J}}{\text{m}^3} \cdot e^{\dfrac{17{,}62 \cdot \vartheta}{\vartheta + 243{,}12}}}
    {461{,}52 \frac{\text{J}}{\text{kg} \cdot \text{K}} \cdot (\vartheta + 273{,}15)\, \text{K}}
    \cdot \frac{10}{1{,}204\, \frac{\text{kg}}{\text{m}^3}}
\end{equation}
mit 
\begin{itemize}
    \item $\vartheta$ : Termodynamische Temperatur in \textdegree C mit $T = (273{,}15 +  \vartheta)\, \text{K}$
    \item Faktor 10: 
    \begin{itemize}
        \item Umrechnung Prozent im Wertebereich $[0 \ldots 100]$ in den Wertebereich $[0 \ldots 1]$
        \item Umrechung der Masse von kg in g 
    \end{itemize}
    \item Die Dichte trockener Luft bei bei Standardbedingungen (20 °C auf Meereshöhe: $101{,}3, \text{Pa}$) beträgt $1,204 \frac{\text{kg}}{\text{m}^3}$. 
\end{itemize}

Die Formel berechnet die absolute Feuchte nicht wie üblich in Gramm pro Kubikmeter ($\mathrm{g/m}^{\mathrm{3}}$), sondern im Mischungsverhältnis Gramm pro Kilogramm trockene Luft ($\mathrm{g/kg}$). Um von der Masse pro Volumen ($\mathrm{m}^{\mathrm{3}}$) auf die Masse pro Luftmasse ($\mathrm{kg}$) zu kommen, muss durch die Dichte der Luft geteilt werden.

Alle Einheiten können gekürzt werden. Es ergibt sich ein Verhältnis der Masse von Wasserdampf zur Masse von trockener Luft. Damit kann die Einheit  $\frac{\text{g}}{\text{kg}}$ gewählt werden.

Damit ist diese Formel gültig für einen Druck von $101{,}3\, \text{Pa}$. Der Druck in der Prozessluft der Dämmschichttrockung wird davon abweichen. Die Größe der Abweichung und deren Einfluss sind zu betrachten.

